\documentclass[12pt]{article}

% Idioma y codificación
\usepackage[spanish, es-tabla, es-notilde]{babel}       %es-tabla para que se titule "Tabla"
\usepackage[utf8]{inputenc}

% Márgenes
\usepackage[a4paper,top=3cm,bottom=2.5cm,left=3cm,right=3cm]{geometry}

% Comentarios de bloque
\usepackage{verbatim}

% Paquetes de links
\usepackage[hidelinks]{hyperref}    % Permite enlaces
\usepackage{url}                    % redirecciona a la web

% Más opciones para enumeraciones
\usepackage{enumitem}

% Personalizar la portada
\usepackage{titling}

% Paquetes de tablas
\usepackage{multirow}

% Para añadir el símbolo de euro
\usepackage{eurosym}

% Paquete para subrayar con diferentes estilos
\usepackage[normalem]{ulem}


%------------------------------------------------------------------------

%Paquetes de figuras
\usepackage{caption}
\usepackage{subcaption} % Figuras al lado de otras
\usepackage{float}      % Poner figuras en el sitio indicado H.


% Paquetes de imágenes
\usepackage{graphicx}       % Paquete para añadir imágenes
\usepackage{transparent}    % Para manejar la opacidad de las figuras

% Paquete para usar colores
\usepackage[dvipsnames, table, xcdraw]{xcolor}
\usepackage{pagecolor}      % Para cambiar el color de la página

% Habilita tamaños de fuente mayores
\usepackage{fix-cm}

% Para los gráficos
\usepackage{tikz}
\usepackage{forest}

% Para poder situar los nodos en los grafos
\usetikzlibrary{positioning}

% Para evitar problemas entre babel y tikz
\usetikzlibrary{babel}


%------------------------------------------------------------------------

% Paquetes de matemáticas
\usepackage{mathtools, amsfonts, amssymb, mathrsfs}
\usepackage[makeroom]{cancel}     % Simplificar tachando
\usepackage{polynom}    % Divisiones y Ruffini
\usepackage{units} % Para poner fracciones diagonales con \nicefrac

\usepackage{pgfplots}   %Representar funciones
\pgfplotsset{compat=1.18}  % Versión 1.18

\usepackage{tikz-cd}    % Para usar diagramas de composiciones
\usetikzlibrary{calc}   % Para usar cálculo de coordenadas en tikz

%Definición de teoremas, etc.
\usepackage{amsthm}
%\swapnumbers   % Intercambia la posición del texto y de la numeración

\theoremstyle{plain}

\makeatletter
\@ifclassloaded{article}{
  % --- Configuración para ARTICLE ---
  \newtheorem{teo}{Teorema}[section]
  \theoremstyle{definition}
  \newtheorem{definicion}{Definición}[section]
  \newtheorem{ejercicio}{Ejercicio}[section]

  \renewcommand{\theteo}{%
    \ifnum\value{section}=0 \arabic{teo}\else \thesection.\arabic{teo}\fi
  }
  \renewcommand{\thedefinicion}{%
    \ifnum\value{section}=0 \arabic{definicion}\else \thesection.\arabic{definicion}\fi
  }
  \renewcommand{\theejercicio}{%
    \ifnum\value{section}=0 \arabic{ejercicio}\else \thesection.\arabic{ejercicio}\fi
  }
}{
  % --- Configuración para BOOK / REPORT ---
  \newtheorem{teo}{Teorema}[chapter]
  \theoremstyle{definition}
  \newtheorem{definicion}{Definición}[chapter]
  \newtheorem{ejercicio}{Ejercicio}[section]

  % Si no hay sección: Capítulo.Teorema (ej. 2.1)
  % Si hay sección:    Capítulo.Sección.Teorema (ej. 2.3.1)
  \renewcommand{\theteo}{%
    \ifnum\value{section}=0
      \thechapter.\arabic{teo}%
    \else
      \thesection.\arabic{teo}% (\thesection ya contiene el número de capítulo)
    \fi
  }
  \renewcommand{\thedefinicion}{%
    \ifnum\value{section}=0
      \thechapter.\arabic{definicion}%
    \else
      \thesection.\arabic{definicion}%
    \fi
  }
  \renewcommand{\theejercicio}{%
    \ifnum\value{section}=0
      \arabic{ejercicio}%
    \else
      \thesection.\arabic{ejercicio}% Formato por defecto: ej. 2.2.1
    \fi
  }
}
\makeatother

% Teoremas dependientes o con contador compartido
\newtheorem{coro}{Corolario}[teo]
\newtheorem{prop}[teo]{Proposición}
\newtheorem{lema}[teo]{Lema}
\newtheorem*{lema*}{Lema}

\theoremstyle{remark}
\newtheorem*{observacion}{Observación}

\theoremstyle{definition}
\newtheorem*{notacion}{Notación}
\newtheorem*{ejemplo}{Ejemplo}
\newtheorem*{ejercicio*}{Ejercicio}


% \renewcommand\qedsymbol{$\blacksquare$}         % Cambiar símbolo QED
%------------------------------------------------------------------------

% Paquetes para encabezados
\usepackage{fancyhdr}
\pagestyle{fancy}
\fancyhf{}

\newcommand{\helv}{ % Modificación tamaño de letra
\fontfamily{}\fontsize{12}{12}\selectfont}
\setlength{\headheight}{15pt} % Amplía el tamaño del índice


%\usepackage{lastpage}   % Referenciar última pag   \pageref{LastPage}
%\fancyfoot[C]{%
%  \begin{minipage}{\textwidth}
%    \centering
%    ~\\
%    \thepage\\
%    \href{https://losdeldgiim.github.io/}{\texttt{\footnotesize losdeldgiim.github.io}}
%  \end{minipage}
%}
\fancyfoot[C]{\thepage}
\fancyfoot[R]{\href{https://losdeldgiim.github.io/}{\texttt{\footnotesize losdeldgiim.github.io}}}

%------------------------------------------------------------------------

% Conseguir que no ponga "Capítulo 1". Sino solo "1."
\makeatletter
\@ifclassloaded{book}{
  \renewcommand{\chaptermark}[1]{\markboth{\thechapter.\ #1}{}} % En el encabezado
    
  \renewcommand{\@makechapterhead}[1]{%
  \vspace*{50\p@}%
  {\parindent \z@ \raggedright \normalfont
    \ifnum \c@secnumdepth >\m@ne
      \huge\bfseries \thechapter.\hspace{1em}\ignorespaces
    \fi
    \interlinepenalty\@M
    \Huge \bfseries #1\par\nobreak
    \vskip 40\p@
  }}
}
\makeatother

%------------------------------------------------------------------------
% Paquetes de cógido
\usepackage{minted}
\renewcommand\listingscaption{Código fuente}

\usepackage{fancyvrb}
% Personaliza el tamaño de los números de línea
\renewcommand{\theFancyVerbLine}{\small\arabic{FancyVerbLine}}

% Estilo para C++
\newminted{cpp}{
    frame=lines,
    framesep=2mm,
    baselinestretch=1.2,
    linenos,
    escapeinside=||
}

% para minted
\definecolor{LightGray}{rgb}{0.95,0.95,0.92}
\setminted{
    linenos=true,
    stepnumber=5,
    numberfirstline=true,
    autogobble,
    breaklines=true,
    breakautoindent=true,
    breaksymbolleft=,
    breaksymbolright=,
    breaksymbolindentleft=0pt,
    breaksymbolindentright=0pt,
    breaksymbolsepleft=0pt,
    breaksymbolsepright=0pt,
    fontsize=\footnotesize,
    bgcolor=LightGray,
    numbersep=10pt
}


\usepackage{listings} % Para incluir código desde un archivo

\renewcommand\lstlistingname{Código Fuente}
\renewcommand\lstlistlistingname{Índice de Códigos Fuente}

% Definir colores
\definecolor{vscodepurple}{rgb}{0.5,0,0.5}
\definecolor{vscodeblue}{rgb}{0,0,0.8}
\definecolor{vscodegreen}{rgb}{0,0.5,0}
\definecolor{vscodegray}{rgb}{0.5,0.5,0.5}
\definecolor{vscodebackground}{rgb}{0.97,0.97,0.97}
\definecolor{vscodelightgray}{rgb}{0.9,0.9,0.9}

% Configuración para el estilo de C similar a VSCode
\lstdefinestyle{vscode_C}{
  backgroundcolor=\color{vscodebackground},
  commentstyle=\color{vscodegreen},
  keywordstyle=\color{vscodeblue},
  numberstyle=\tiny\color{vscodegray},
  stringstyle=\color{vscodepurple},
  basicstyle=\scriptsize\ttfamily,
  breakatwhitespace=false,
  breaklines=true,
  captionpos=b,
  keepspaces=true,
  numbers=left,
  numbersep=5pt,
  showspaces=false,
  showstringspaces=false,
  showtabs=false,
  tabsize=2,
  frame=tb,
  framerule=0pt,
  aboveskip=10pt,
  belowskip=10pt,
  xleftmargin=10pt,
  xrightmargin=10pt,
  framexleftmargin=10pt,
  framexrightmargin=10pt,
  framesep=0pt,
  rulecolor=\color{vscodelightgray},
  backgroundcolor=\color{vscodebackground},
}

%------------------------------------------------------------------------
% La etiqueda IA sirve para indicar a partir de donde le pasamos el preambulo a la IA para el escaneo
% IA

% Comandos definidos
\newcommand{\bb}[1]{\mathbb{#1}}
\newcommand{\cc}[1]{\mathcal{#1}}

% I prefer the slanted \leq
\let\oldleq\leq % save them in case they're every wanted
\let\oldgeq\geq
\renewcommand{\leq}{\leqslant}
\renewcommand{\geq}{\geqslant}

% Si y solo si
\newcommand{\sii}{\iff}

% MCD y MCM
\DeclareMathOperator{\mcd}{mcd}
\DeclareMathOperator{\mcm}{mcm}

% Signo
\DeclareMathOperator{\sgn}{sgn}

% Imagen
\DeclareMathOperator{\Ima}{Im}

% Aplicación identidad
\newcommand{\id}{\mathrm{Id}}

% Letras griegas
\newcommand{\eps}{\epsilon}
\newcommand{\veps}{\varepsilon}
\newcommand{\lm}{\lambda}

\newcommand{\ol}{\overline}
\newcommand{\ul}{\underline}
\newcommand{\wt}{\widetilde}
\newcommand{\wh}{\widehat}

\let\oldvec\vec
\renewcommand{\vec}{\overrightarrow}

% Derivadas parciales
\newcommand{\del}[2]{\frac{\partial #1}{\partial #2}}
\newcommand{\Del}[3]{\frac{\partial^{#1} #2}{\partial #3^{#1}}}
\newcommand{\deld}[2]{\dfrac{\partial #1}{\partial #2}}
\newcommand{\Deld}[3]{\dfrac{\partial^{#1} #2}{\partial #3^{#1}}}


\newcommand{\AstIg}{\stackrel{(\ast)}{=}}
\newcommand{\Hop}{\stackrel{L'H\hat{o}pital}{=}}

\newcommand{\red}[1]{{\color{red}#1}} % Para integrales, destacar los cambios.

% Método de integración
\newcommand{\MetInt}[2]{
    \left[\begin{array}{c}
        #1 \\ #2
    \end{array}\right]
}

% Declarar aplicaciones
% 1. Nombre aplicación
% 2. Dominio
% 3. Codominio
% 4. Variable
% 5. Imagen de la variable
\newcommand{\Func}[5]{
    \begin{equation*}
        \begin{array}{rrll}
            \displaystyle #1:& \displaystyle  #2 & \longrightarrow & \displaystyle  #3\\
               & \displaystyle  #4 & \longmapsto & \displaystyle  #5
        \end{array}
    \end{equation*}
}

%------------------------------------------------------------------------


\newcommand{\myhref}[2]{\href{#1}{\color{blue}\uline{#2}}}


\begin{document}

    % 1. Foto de fondo
    % 2. Título
    % 3. Encabezado Izquierdo
    % 4. Color de fondo
    % 5. Coord x del titulo
    % 6. Coord y del titulo
    % 7. Fecha

    
    % 1. Foto de fondo
% 2. Título
% 3. Encabezado Izquierdo
% 4. Color de fondo
% 5. Coord x del titulo
% 6. Coord y del titulo
% 7. Fecha
% 8. Autor

\newcommand{\portada}[8]{

    \portadaBase{#1}{#2}{#3}{#4}{#5}{#6}{#7}{#8}
    \portadaBook{#1}{#2}{#3}{#4}{#5}{#6}{#7}{#8}
}

\newcommand{\portadaExamen}[8]{

    \portadaBase{#1}{#2}{#3}{#4}{#5}{#6}{#7}{#8}
    \portadaArticle{#1}{#2}{#3}{#4}{#5}{#6}{#7}{#8}
}




\newcommand{\portadaBase}[8]{

    % Tiene la portada principal y la licencia Creative Commons
    
    % 1. Foto de fondo
    % 2. Título
    % 3. Encabezado Izquierdo
    % 4. Color de fondo
    % 5. Coord x del titulo
    % 6. Coord y del titulo
    % 7. Fecha
    % 8. Autor
    
    
    \thispagestyle{empty}               % Sin encabezado ni pie de página
    \newgeometry{margin=0cm}        % Márgenes nulos para la primera página
    
    
    % Encabezado
    \fancyhead[L]{\helv #3}
    \fancyhead[R]{\helv \nouppercase{\leftmark}}
    
    
    \pagecolor{#4}        % Color de fondo para la portada
    
    \begin{figure}[p]
        \centering
        \transparent{0.3}           % Opacidad del 30% para la imagen
        
        \includegraphics[width=\paperwidth, keepaspectratio]{../../_assets/#1}
    
        \begin{tikzpicture}[remember picture, overlay]
            \node[anchor=north west, text=white, opacity=1, font=\fontsize{60}{90}\selectfont\bfseries\sffamily, align=left] at (#5, #6) {#2};
            
            \node[anchor=south east, text=white, opacity=1, font=\fontsize{12}{18}\selectfont\sffamily, align=right] at (9.7, 3) {\href{https://losdeldgiim.github.io/}{\textbf{Los Del DGIIM}, \texttt{\footnotesize losdeldgiim.github.io}}};
            
            \node[anchor=south east, text=white, opacity=1, font=\fontsize{12}{15}\selectfont\sffamily, align=right] at (9.7, 1.8) {Doble Grado en Ingeniería Informática y Matemáticas\\Universidad de Granada};
        \end{tikzpicture}
    \end{figure}
    
    
    \restoregeometry        % Restaurar márgenes normales para las páginas subsiguientes
    \pagecolor{white}       % Restaurar el color de página
    
    
    \newpage
    \thispagestyle{empty}               % Sin encabezado ni pie de página
    \begin{tikzpicture}[remember picture, overlay]
        \node[anchor=south west, inner sep=3cm] at (current page.south west) {
            \begin{minipage}{0.5\paperwidth}
                \href{https://creativecommons.org/licenses/by-nc-nd/4.0/}{
                    \includegraphics[height=2cm]{../../_assets/Licencia.png}
                }\vspace{1cm}\\
                Esta obra está bajo una
                \href{https://creativecommons.org/licenses/by-nc-nd/4.0/}{
                    Licencia Creative Commons Atribución-NoComercial-SinDerivadas 4.0 Internacional (CC BY-NC-ND 4.0).
                }\\
    
                Eres libre de compartir y redistribuir el contenido de esta obra en cualquier medio o formato, siempre y cuando des el crédito adecuado a los autores originales y no persigas fines comerciales. 
            \end{minipage}
        };
    \end{tikzpicture}
    
    
    
    % 1. Foto de fondo
    % 2. Título
    % 3. Encabezado Izquierdo
    % 4. Color de fondo
    % 5. Coord x del titulo
    % 6. Coord y del titulo
    % 7. Fecha
    % 8. Autor


}


\newcommand{\portadaBook}[8]{

    % 1. Foto de fondo
    % 2. Título
    % 3. Encabezado Izquierdo
    % 4. Color de fondo
    % 5. Coord x del titulo
    % 6. Coord y del titulo
    % 7. Fecha
    % 8. Autor

    % Personaliza el formato del título
    \pretitle{\begin{center}\bfseries\fontsize{42}{56}\selectfont}
    \posttitle{\par\end{center}\vspace{2em}}
    
    % Personaliza el formato del autor
    \preauthor{\begin{center}\Large}
    \postauthor{\par\end{center}\vfill}
    
    % Personaliza el formato de la fecha
    \predate{\begin{center}\huge}
    \postdate{\par\end{center}\vspace{2em}}
    
    \title{#2}
    \author{\href{https://losdeldgiim.github.io/}{Los Del DGIIM, \texttt{\large losdeldgiim.github.io}}
    \\ \vspace{0.5cm}#8}
    \date{Granada, #7}
    \maketitle
    
    \tableofcontents
}




\newcommand{\portadaArticle}[8]{

    % 1. Foto de fondo
    % 2. Título
    % 3. Encabezado Izquierdo
    % 4. Color de fondo
    % 5. Coord x del titulo
    % 6. Coord y del titulo
    % 7. Fecha
    % 8. Autor

    % Personaliza el formato del título
    \pretitle{\begin{center}\bfseries\fontsize{42}{56}\selectfont}
    \posttitle{\par\end{center}\vspace{2em}}
    
    % Personaliza el formato del autor
    \preauthor{\begin{center}\Large}
    \postauthor{\par\end{center}\vspace{3em}}
    
    % Personaliza el formato de la fecha
    \predate{\begin{center}\huge}
    \postdate{\par\end{center}\vspace{5em}}
    
    \title{#2}
    \author{\href{https://losdeldgiim.github.io/}{Los Del DGIIM, \texttt{\large losdeldgiim.github.io}}
    \\ \vspace{0.5cm}#8}
    \date{Granada, #7}
    \thispagestyle{empty}               % Sin encabezado ni pie de página
    \maketitle
    \vfill
}
    \portadaExamen{ffccA4.jpg}{Análisis Funcional\\Examen IV}{Análisis Funcional. Examen IV}{MidnightBlue}{-8}{28}{2026}{Daniel Morán Sánchez}

    \begin{description}
        \item[Asignatura] Análisis Funcional.
        \item[Curso Académico] 2024/25.
        \item[Grado] Doble Grado en Ingeniería Informática y Matemáticas.
        % \item[Grupo] ---.
        % \item[Profesor] ---.
        \item[Descripción] Examen Ordinario.
        % \item[Fecha] 15 de enero de 2025.
        % \item[Duración] ---.
    
    \end{description}
    \newpage


    % ------------------------------------
    
    \begin{ejercicio}[3 puntos]
        Sea $X = C([0,1],\mathbb{R})$ el espacio vectorial de las funciones continuas de $[0,1]$ a $\mathbb{R}$. Se define
        \begin{equation*}
            \|f\|_1 := \int_{0}^{1} |f(t)|~dt , \qquad \|f\|_\infty := \sup\{|f(t)| : t\in [0,1]\}\qquad \forall f\in X
        \end{equation*}
        y se escribe $X_1 = (X,\|\cdot \|_1)$, $X_{\infty} = (X,\|\cdot \|_\infty)$. Definimos el funcional lineal $\delta:X\to \mathbb{R}$ por
        \begin{equation*}
            \delta(f) = f(0) \qquad \forall f\in X
        \end{equation*}
        \begin{enumerate}[label=\alph*)]
            \item \textbf{[0.75 puntos]} ¿Son equivalentes $\|\cdot \|_1$ y $\|\cdot \|_\infty$ en $X$? Justifica tu respuesta.
            \item \textbf{[0.75 puntos]} ¿Es $X_1$ completo? Justifica tu respuesta.
            \item \textbf{[0.75 puntos]} Demuestra que $\delta\in X_\infty^\ast$ y calcula su norma.
            \item \textbf{[0.75 puntos]} ¿Es $\delta$ continuo respecto a $\|\cdot \|_1$? Justifica tu respuesta.
        \end{enumerate}
    \end{ejercicio}

    \begin{ejercicio}[3 puntos]
        Sean $S:c_0\to c_0$ y $T:l_1\to l_1$ operadores lineales y supongamos que
        \begin{equation*}
            \sum_{k=1}^{\infty}[Sx](k)y(k) = \sum_{k=1}^{\infty}x(k)[Ty](k) \qquad \forall x\in c_0, \quad \forall y\in l_1
        \end{equation*}
        Demuestra que $S$ y $T$ son continuos.
    \end{ejercicio}

    \begin{ejercicio}[4 puntos]
        Sean $1<p<\infty$ ($p' = \nicefrac{p}{p-1}$), $\emptyset \neq \Omega\subset \mathbb{R}^N$ un subconjunto medible y $T:L^{p'}(\Omega)\to L^p(\Omega)^\ast$ definida mediante
        \begin{equation*}
            \langle Tu,f \rangle  = \int_\Omega uf \qquad \forall u\in L^{p'}(\Omega), \quad \forall f\in L^p(\Omega)
        \end{equation*}
        \begin{enumerate}[label=\alph*)]
            \item \textbf{[1 punto]} Prueba que $\|Tu\|_{L^p(\Omega)^\ast} = \|u\|_{L^{p'}(\Omega)}$ y que $T$ es inyectivo.
            \item \textbf{[1 punto]} Usando que $L^p(\Omega)$ es reflexivo, prueba que $T(L^{p'}(\Omega)) = L^p(\Omega)^\ast$.
        \end{enumerate}
        Si para $\Omega = \left]0,1\right[$ consideramos la sucesión de funciones $\{g_n\}$ definidas mediante $g_n(x) = n^{\nicefrac{1}{p}}e^{-nx}$ para todo $x\in \left]0,1\right[$:
        \begin{enumerate}[label=\alph*)]
            \item[c)] \textbf{[0.5 puntos]} prueba que $\{g_n\}\to 0$ a.e. y que $\{g_n\}$ es acotada en $L^p(\Omega)$.
            \item[d)] \textbf{[0.5 puntos]} Demuestra que $\{g_n\}\not\to 0$ en $L^p(\Omega)$.
            \item[e)] \textbf{[1 punto]} Usando que $C_c(0,1)$ es denso en $L^{p'}(\Omega)$, prueba que $\{g_n\}\to 0$ en la topología débil $\sigma(L^p(\Omega),L^{p'}(\Omega))$.
        \end{enumerate}
    \end{ejercicio}


    \newpage
    \setcounter{ejercicio}{0} % Reiniciar contador de ejercicios

    \begin{ejercicio}[3 puntos]
        Sea $X = C([0,1],\mathbb{R})$ el espacio vectorial de las funciones continuas de $[0,1]$ a $\mathbb{R}$. Se define
        \begin{equation*}
            \|f\|_1 := \int_{0}^{1} |f(t)|~dt , \qquad \|f\|_\infty := \sup\{|f(t)| : t\in [0,1]\}\qquad \forall f\in X
        \end{equation*}
        y se escribe $X_1 = (X,\|\cdot \|_1)$, $X_{\infty} = (X,\|\cdot \|_\infty)$. Definimos el funcional lineal $\delta:X\to \mathbb{R}$ por
        \begin{equation*}
            \delta(f) = f(0) \qquad \forall f\in X
        \end{equation*}
        \begin{enumerate}[label=\alph*)]
            \item \textbf{[0.75 puntos]} ¿Son equivalentes $\|\cdot \|_1$ y $\|\cdot \|_\infty$ en $X$? Justifica tu respuesta.
            \item \textbf{[0.75 puntos]} ¿Es $X_1$ completo? Justifica tu respuesta.
            \item \textbf{[0.75 puntos]} Demuestra que $\delta\in X_\infty^\ast$ y calcula su norma.
            \item \textbf{[0.75 puntos]} ¿Es $\delta$ continuo respecto a $\|\cdot \|_1$? Justifica tu respuesta.
        \end{enumerate}


        Ver resolución del ejercicio 1 del Examen 3.
    \end{ejercicio}
    \newpage

    \begin{ejercicio}[3 puntos]
        Sean $S:c_0\to c_0$ y $T:l_1\to l_1$ operadores lineales y supongamos que
        \begin{equation*}
            \sum_{k=1}^{\infty}[Sx](k)y(k) = \sum_{k=1}^{\infty}x(k)[Ty](k) \qquad \forall x\in c_0, \quad \forall y\in l_1
        \end{equation*}
        Demuestra que $S$ y $T$ son continuos.\\

        Ver resolución del ejercicio 3 del Examen 3.
    \end{ejercicio}

    \newpage
    \begin{ejercicio}[4 puntos]
        Sean $1<p<\infty$ ($p'=\nicefrac{p}{p-1}$), $\emptyset\neq\Omega\subset\mathbb{R}^N$ un conjunto medible y
        $T:L^{p'}(\Omega)\to L^p(\Omega)^\ast$ definido por
        \begin{equation*}
        \langle Tu,f\rangle := \int_\Omega u f,
        \qquad
        \forall u\in L^{p'}(\Omega),\ \forall f\in L^p(\Omega).
        \end{equation*}
        
        \begin{enumerate}[label=\alph*)]
        
        \item \textbf{Prueba que $\|Tu\|_{L^p(\Omega)^\ast}=\|u\|_{L^{p'}(\Omega)}$ y que $T$ es inyectivo.}
        
        Sea $u\in L^{p'}(\Omega)$. Por definición de la norma en el dual,
        \begin{equation*}
        \|Tu\|_{(L^p(\Omega))^\ast}
        =
        \sup_{\|f\|_{L^p(\Omega)}\le 1}
        \left|\int_\Omega u f\right|.
        \end{equation*}
        Por la desigualdad de Hölder,
        \begin{equation*}
        \left|\int_\Omega u f\right|
        \le
        \|u\|_{L^{p'}(\Omega)}\,\|f\|_{L^p(\Omega)},
        \end{equation*}
        y por tanto
        \begin{equation*}
        \|Tu\|_{(L^p(\Omega))^\ast}\le \|u\|_{L^{p'}(\Omega)}.
        \end{equation*}
        
        Veamos ahora la desigualdad contraria. Si $u=0$ es trivial. Supongamos $u\neq 0$ y definimos
        \begin{equation*}
        f_0:=u\,|u|^{p'-2}=\operatorname{sgn}(u)\,|u|^{p'-1}.
        \end{equation*}
        Primero probamos que $f_0\in L^p(\Omega)$. En efecto,
        \begin{equation*}
        |f_0|^p = |u|^{p(p'-1)}.
        \end{equation*}
        Como $p'=\frac{p}{p-1}$, se tiene $p'-1=\frac{1}{p-1}$ y por tanto
        \begin{equation*}
        p(p'-1)=\frac{p}{p-1}=p'.
        \end{equation*}
        Luego
        \begin{equation*}
        \int_\Omega |f_0|^p=\int_\Omega |u|^{p'}<\infty,
        \end{equation*}
        y así $f_0\in L^p(\Omega)$. Además,
        \begin{equation*}
        \|f_0\|_{L^p(\Omega)}^p=\int_\Omega |u|^{p'}=\|u\|_{L^{p'}(\Omega)}^{p'}
        \quad\Longrightarrow\quad
        \|f_0\|_{L^p(\Omega)}=\|u\|_{L^{p'}(\Omega)}^{\frac{p'}{p}}
        =\|u\|_{L^{p'}(\Omega)}^{p'-1}.
        \end{equation*}
        
        Ahora evaluamos el funcional en $f_0$:
        \begin{equation*}
        \int_\Omega u f_0=\int_\Omega u\,u\,|u|^{p'-2}=\int_\Omega |u|^{p'}
        =\|u\|_{L^{p'}(\Omega)}^{p'}.
        \end{equation*}
        Consideramos entonces $\displaystyle f:=\frac{f_0}{\|f_0\|_{L^p(\Omega)}}$, que satisface $\|f\|_{L^p(\Omega)}=1$. Se tiene:
        \begin{multline*}
        \|Tu\|_{(L^p(\Omega))^\ast}
        =
        \sup_{\|g\|_{L^p(\Omega)}\le 1}
        \left|\langle Tu,g\rangle\right|
        =
        \sup_{\|g\|_{L^p(\Omega)}\le 1}
        \left|\int_\Omega u g\right| \\
        \ge
        \left|\int_\Omega u\, f\right|
        =
        \frac{\int_\Omega |u|^{p'}}{\|f_0\|_{L^p(\Omega)}}
        =
        \frac{\|u\|_{L^{p'}(\Omega)}^{p'}}{\|u\|_{L^{p'}(\Omega)}^{p'-1}}
        =
        \|u\|_{L^{p'}(\Omega)}.
        \end{multline*}
        
        Junto con la desigualdad anterior concluimos
        \begin{equation*}
        \|Tu\|_{(L^p(\Omega))^\ast}=\|u\|_{L^{p'}(\Omega)}.
        \end{equation*}
        
        Finalmente, $T$ es una isometría lineal, y sabemos que toda isometría lineal es inyectiva: si $Tu=0$ porque
        \begin{equation*}
        0=\|Tu\|_{(L^p(\Omega))^\ast}=\|u\|_{L^{p'}(\Omega)}
        \quad\Longrightarrow\quad
        u=0 \ \text{a.e.}
        \end{equation*}
        %-------------------------------------------------
        \item \textbf{Usando que $L^p(\Omega)$ es reflexivo, prueba que $T(L^{p'}(\Omega))=L^p(\Omega)^\ast$.}
        
        Por el apartado anterior, $T$ es una isometría lineal. Como $L^{p'}(\Omega)$ es Banach,
        se sigue que $T(L^{p'}(\Omega))$ es Banach y por tanto un subespacio cerrado de $L^p(\Omega)^\ast$.\\
        Para ver que es todo el dual, basta probar que $T(L^{p'}(\Omega))$ es denso en $L^p(\Omega)^\ast$.
        
        Usamos el último corolario del teorema de Hahn--Banach: un subespacio $M$ de un espacio normado $E$ es denso si y solo si el único funcional continuo de $E^\ast$ que se anula en $M$ es el funcional nulo.


        Sea
        $\Phi\in L^p(\Omega)^{\ast\ast}$ tal que
        \begin{equation*}
        \langle \Phi, Tu\rangle = 0
        \qquad \forall u\in L^{p'}(\Omega).
        \end{equation*}
        Como $L^p(\Omega)$ es reflexivo, existe $f\in L^p(\Omega)$ tal que
        $\Phi = J(f)$, donde $J$ es la inmersión canónica. Entonces
        \begin{equation*}
        0
        =
        \langle J(f),Tu\rangle
        =
        \langle Tu,f\rangle
        =
        \int_\Omega u f
        \qquad \forall u\in L^{p'}(\Omega).
        \end{equation*}
        Esto implica que $f=0$ $a.e.$ , luego $\Phi=0$.
        Por el teorema de Hahn--Banach, $T(L^{p'}(\Omega))$ es denso pero además sabíamos que era cerrado,
        por lo que
        \begin{equation*}
        T(L^{p'}(\Omega))=L^p(\Omega)^\ast.
        \end{equation*}
        
        \end{enumerate}
        
        \vspace{2mm}
        
        Ahora supongamos $\Omega=(0,1)$ y definamos
        \begin{equation*}
        g_n(x)=n^{1/p}e^{-nx}, \qquad x\in(0,1).
        \end{equation*}
        
        \begin{enumerate}[label=\alph*),resume]
        
        %-------------------------------------------------
        \item \textbf{Prueba que $g_n\to 0$ a.e. y que $\{g_n\}$ es acotada en $L^p(\Omega)$.}
        
        Para todo $x\in(0,1)$,
        \begin{equation*}
        g_n(x)=n^{1/p}e^{-nx}\xrightarrow[n\to\infty]{}0,
        \end{equation*}
        ya que la exponencial domina a cualquier potencia. Luego $g_n\to 0$ a.e.
        
        Además,
        \begin{equation*}
        \|g_n\|_{L^p(\Omega)}^p
        =
        \int_0^1 n e^{-npx}\,dx
        =
        \frac{1}{p}\left(1-e^{-np}\right)
        \le
        \frac{1}{p}.
        \end{equation*}
        Por tanto, $\{g_n\}$ es acotada en $L^p(\Omega)$.
        
        %-------------------------------------------------
        \item \textbf{Demuestra que $g_n\not\to 0$ en $L^p(\Omega)$.}
        
        Se tiene
        \begin{equation*}
        \|g_n\|_{L^p(\Omega)}
        =
        \left(\frac{1}{p}(1-e^{-np})\right)^{1/p}
        \xrightarrow[n\to\infty]{}
        \frac{1}{p^{1/p}}\neq 0.
        \end{equation*}
        Luego $\|g_n\|_{L^p(\Omega)}$ no converge a $0$, por tanto .
        
        %-------------------------------------------------
        \item \textbf{Usando que $C_c(0,1)$ es denso en $L^{p'}(\Omega)$, prueba que $g_n\to 0$
        en la topología débil $\sigma(L^p(\Omega),L^{p'}(\Omega))$.}

        Recordemos que $g_n \rightharpoonup 0$ en $\sigma(L^p(\Omega),L^{p'}(\Omega))$ significa que
        \[
        \int_0^1 g_n(x)\,u(x)\,dx \longrightarrow 0
        \qquad \forall\,u\in L^{p'}(0,1),
        \]
        es decir, que toda forma lineal continua sobre $L^p$ (identificada con un $u\in L^{p'}$) aplicada a $g_n$
        converge a $0$.
        
        Sea $u\in L^{p'}(\Omega)$ y $\varepsilon>0$. Como $C_c(0,1)$ es denso en $L^{p'}(\Omega)$,
        existe $\varphi\in C_c(0,1)$ tal que
        \begin{equation*}
        \|u-\varphi\|_{L^{p'}(\Omega)}<\varepsilon.
        \end{equation*}
        Entonces,
        \begin{equation*}
        \left|\int_0^1 g_n u\right|
        \le
        \left|\int_0^1 g_n (u-\varphi)\right|
        +
        \left|\int_0^1 g_n \varphi\right|.
        \end{equation*}
        Para el primer término, por Hölder,
        \begin{equation*}
        \left|\int_0^1 g_n (u-\varphi)\right|
        \le
        \|g_n\|_{L^p(\Omega)}\|u-\varphi\|_{L^{p'}(\Omega)}
        \le
        \frac{1}{p^{1/p}}\varepsilon \qquad \forall \varepsilon >0, \forall u \in L^{p'}(\Omega)
        \end{equation*}
        Para el segundo término, como $\varphi$ es continua con soporte compacto y $g_n\to 0$ a.e.,
        por convergencia dominada,
        \begin{equation*}
        \int_0^1 g_n \varphi \xrightarrow[n\to\infty]{} 0.
        \end{equation*}
        Concluimos que
        \begin{equation*}
        \lim_{n\to\infty}\int_0^1 g_n u = 0
        \qquad \forall u\in L^{p'}(\Omega),
        \end{equation*}        
        es decir,
        \begin{equation*}
        g_n \rightharpoonup 0
        \quad \text{en } \sigma(L^p(\Omega),L^{p'}(\Omega)).
        \end{equation*}
        \end{enumerate}
    \end{ejercicio}
    

\end{document}
